% Minimal preamble for Open Logic Project sources (pinned commit 1e960beff9ed7835bf3e3f1335e21af3439cd107).
%
% OLP's own master preamble cannot run on the project host: TeX Live 2025 there lacks
% prooftree.sty, ntheorem.sty, extarrows.sty, mathdots.sty, isomath.sty and subfiles.sty.
% Following the project's own rule ("if the author preamble needs packages the toolchain lacks,
% do not invent macros -- use a minimal working preamble"), this file declares, verbatim in
% spirit and numbering behaviour, exactly the theorem-like environments OLP defines in
% sty/open-logic-book-envs.sty (thm/lem/prop/cor/defn/ex/prob/rem/note/case/axiom/conv),
% plus the shorthand accents OLP uses (`!!a`, `!!b`, `!!{...}`).
%
% Author text is untouched; only the packaging is ours.
\documentclass[10pt,a4paper]{article}

\usepackage{amsmath}
\usepackage{amssymb}
\usepackage{amsthm}
\usepackage{graphicx}
\usepackage{hyperref}
\usepackage{enumitem}
\usepackage[framemethod=default]{mdframed}

% ---- OLP shorthand accents (from sty/open-logic.sty) ----
% NB: \! is taken by LaTeX (negative thin space), so we only define the !!a / !!b families
% and the \olup helper used by them.
\newcommand*{\olup}[1]{\MakeUppercase{#1}}
\def\!!!a#1{\olup{#1}}
\def\!a#1{\olup{#1}}
\def\!b#1{\olup{#1}}

% ---- theorem-like environments, same family numbering as OLP ----
\theoremstyle{plain}
\newtheorem{thm}{Theorem}[section]
\newtheorem{lem}[thm]{Lemma}
\newtheorem{prop}[thm]{Proposition}
\newtheorem{cor}[thm]{Corollary}
\newtheorem{axiom}{Axiom}

\theoremstyle{definition}
\newtheorem{defn}[thm]{Definition}
\newtheorem{ex}[thm]{Example}
\newtheorem{prob}{Problem}[section]
\newtheorem{probd}{Problem}
\newtheorem{conv}[thm]{Convention}

\theoremstyle{remark}
\newtheorem{rem}[thm]{Remark}
\newtheorem{note}[thm]{Note}
\newtheorem{case}[thm]{Case}

% ---- environments OLP defines as wrappers ----
\newenvironment{probtag}[1]{\begin{prob}}{\end{prob}}
\newenvironment{explain}{}{}
\newenvironment{explain*}{}{}
\newenvironment{editorial}{}{}
\newenvironment{digress}{}{}
\newenvironment{digress*}{}{}
\newenvironment{history}{}{}
\newenvironment{reference}{}{}
\newcommand*{\olfileid}[3]{}
\newcommand*{\olsection}[1]{\section*{#1}}
\newcommand*{\olchapter}[3]{}
\newcommand*{\olpart}[2]{}
\newcommand*{\olimport}[1]{}
\newcommand*{\olinclude}[1]{}
\newcommand*{\olpath}[1]{}
\newcommand*{\ollangid}[1]{}
\newcommand*{\problemtag}[1]{}
\newcommand*{\problemauthor}[1]{}
\newcommand*{\ollabel}[1]{\label{#1}}
% \olref takes 0-3 OPTIONAL arguments plus one mandatory label:
%     \olref{lab}  |  \olref[from]{lab}  |  \olref[from][to]{lab}  |  \olref[a][b][c]{lab}
% A single-argument definition silently swallows "[sth]" and emits a broken reference
% (measured on the host: the text came out as "??iso]x"). The chain below consumes the
% brackets and cross-references only the mandatory argument.
\makeatletter
\newcommand{\olref}{\@ifnextchar[{\olref@i}{\olref@final}}
\def\olref@i[#1]{\@ifnextchar[{\olref@ii}{\olref@final}}
\def\olref@ii[#1]{\@ifnextchar[{\olref@iii}{\olref@final}}
\def\olref@iii[#1]{\olref@final}
\def\olref@final#1{\ref{#1}}
\makeatother
\newcommand*{\olnewpage}{\clearpage}

% ===================== OLP compatibility layer (1) =====================

\makeatletter

% ---------------- proof trees (OLP's prooftree is a bussproofs derivative) ----------------
\providecommand{\AxiomC}[1]{}
\providecommand{\UnaryInfC}[1]{}
\providecommand{\BinaryInfC}[1]{}
\providecommand{\TrinaryInfC}[1]{}
\providecommand{\QuaternaryInfC}[1]{}
\providecommand{\QuinaryInfC}[1]{}
\providecommand{\RightLabel}[1]{}
\providecommand{\LeftLabel}[1]{}
\providecommand{\doubleLine}{}
\providecommand{\noLine}{}
\providecommand{\DisplayProof}{}
\providecommand{\fCenter}{}
\newenvironment{prooftree}{}{}
\newenvironment{ltprooftree}{}{}
\newenvironment{oltableau}{}{}

% ---------------- logical notation ----------------
\providecommand{\lif}{\mathbin{\rightarrow}}
\providecommand{\lforall}{\mathord{\forall}}
\providecommand{\lexists}{\mathord{\exists}}
\providecommand{\ltrue}{\top}
\providecommand{\lfalse}{\bot}
\providecommand{\Ltrue}{\top}
\providecommand{\Lfalse}{\bot}
\providecommand{\Sequent}{\mathrel{\Rightarrow}}
\providecommand{\Proves}{\mathrel{\vdash}}
\providecommand{\Entails}{\mathrel{\models}}
\providecommand{\Subst}[3]{#1[#2/#3]}
\providecommand{\sFmla}[2]{#1(#2)}
\providecommand{\pFmla}[2]{#1(#2)}
\providecommand{\formula}{\relax}
\providecommand{\True}{\top}
\providecommand{\False}{\bot}
\providecommand{\Args}[1]{#1}
\providecommand{\Ts}{\ensuremath{\top}}
\providecommand{\Fi}{\ensuremath{\bot}}
\providecommand{\lfalsum}{\bot}

% ---------------- objects, structures, metatheory ----------------
\providecommand{\Obj}{}
\providecommand{\Struct}{\ensuremath{\mathfrak{A}}}
\providecommand{\Domain}{\ensuremath{|\mathfrak{A}|}}
\providecommand{\Assign}{\ensuremath{s}}
\providecommand{\Sat}[2]{\ensuremath{#1 \models #2}}
\providecommand{\mSat}[3]{\ensuremath{#1 \models #2[#3]}}
\providecommand{\mModel}[2]{\ensuremath{#1 \models #2}}
\providecommand{\Lang}{\ensuremath{\mathcal{L}}}
\providecommand{\Th}{\ensuremath{\mathrm{Th}}}
\providecommand{\Log}{\ensuremath{\mathrm{L}}}
\providecommand{\tuple}[1]{\langle #1 \rangle}
\providecommand{\num}[1]{\underline{#1}}
\providecommand{\cfind}[1]{\varphi_{#1}}
\providecommand{\Nat}{\ensuremath{\mathbb{N}}}
\providecommand{\Mult}{\ensuremath{\mathrm{Mult}}}
\providecommand{\Add}{\ensuremath{\mathrm{Add}}}
\providecommand{\Exp}{\ensuremath{\mathrm{Exp}}}
\providecommand{\Zero}{\ensuremath{0}}
\providecommand{\Succ}{\ensuremath{\mathrm{S}}}
\providecommand{\lambd}{\ensuremath{\lambda}}
\providecommand{\bredpar}{\mathrel{\rightarrow_\beta}}
\providecommand{\bcd}[1]{\ensuremath{\ulcorner #1 \urcorner}}
\providecommand{\Set}[1]{\ensuremath{\{#1\}}}
\providecommand{\indcase}[1]{\par\noindent\emph{Case #1}.\quad}
\providecommand{\indfrm}[1]{\par\noindent\emph{Form #1}.\quad}
\providecommand{\iftag}[1]{#1}
\providecommand{\tagitem}[1]{\item[#1]}
\providecommand{\Elim}[1]{}
\providecommand{\Intro}[1]{}
\providecommand{\LeftR}[1]{}
\providecommand{\RightR}[1]{}
\providecommand{\TAss}[1]{}
\providecommand{\TRule}[1]{}

\makeatother
\providecommand{\olpath}{.}

% ===================== OLP compatibility layer (2) =====================
% --- second compatibility block: notation and shorthands used by the frozen OLP units ---
\providecommand{\CutCS}{\mathrm{Cut}}
\providecommand{\LeftSubproofLabel}{\relax}
\providecommand{\RightSubproofLabel}{\relax}
\providecommand{\maeh}{\relax}
\providecommand{\cheight}{\relax}
\providecommand{\ande}{\mathrm{and}_E}
\providecommand{\andi}{\mathrm{and}_I}
\providecommand{\ore}{\mathrm{or}_E}
\providecommand{\ori}{\mathrm{or}_I}
\providecommand{\typeof}{\mathord{:}}
\providecommand{\pValue}[1]{#1}
\providecommand{\Sent}{\mathrm{Sent}}
\providecommand{\cardplus}{\mathbin{+}}
\providecommand{\cardtimes}{\mathbin{\cdot}}
\providecommand{\cardexpo}{\mathbin{\hat{\ }}}
\providecommand{\cardnless}{\not<}
\providecommand{\cardle}{\leq}
\providecommand{\funfromto}[3]{#1\colon #2 \to #3}
\providecommand{\funimage}[1]{\operatorname{im}(#1)}
\providecommand{\fdefined}{\mathord{\downarrow}}
\providecommand{\canonord}{\operatorname{ord}}
\providecommand{\ordtimes}{\mathbin{\cdot}}
\providecommand{\ordeq}[2]{#1 = #2}
\providecommand{\ordneq}[2]{#1 \neq #2}
\providecommand{\ordtype}[1]{\operatorname{ord}(#1)}
\providecommand{\ordsucc}{\mathord{+1}}
\providecommand{\ordunion}{\cup}
\providecommand{\onesphere}{S^1}
\providecommand{\rotationsgroup}{\mathrm{SO}(2)}
\providecommand{\equivrep}[1]{[#1]}
\providecommand{\rlexless}{\prec}
\providecommand{\ZFC}{\mathrm{ZFC}}
\providecommand{\ZF}{\mathrm{ZF}}
\providecommand{\Rat}{\mathbb{Q}}
\providecommand{\Real}{\mathbb{R}}
\providecommand{\PosInt}{\mathbb{Z}^{+}}
\providecommand{\Id}{\mathrm{id}}
\providecommand{\Bin}{\mathrm{Bin}}
\providecommand{\pto}{\rightharpoonup}
\providecommand{\dom}{\operatorname{dom}}
\providecommand{\ran}{\operatorname{ran}}
\providecommand{\len}{\operatorname{len}}
\providecommand{\concat}{\mathbin{{}^\frown}}
\providecommand{\card}[1]{\lvert #1 \rvert}
\providecommand{\Pow}{\mathcal{P}}
\providecommand{\Setabs}[2]{\{\, #1 \mid #2 \,\}}
\providecommand{\comp}{\circ}
\providecommand{\Char}{\chi}
\providecommand{\Prf}{\operatorname{Prf}}
\providecommand{\Prov}{\operatorname{Prov}}
\providecommand{\Proj}{\operatorname{proj}}
\providecommand{\Gn}{\operatorname{gn}}
\providecommand{\gn}{\operatorname{gn}}
\providecommand{\Value}[1]{\operatorname{val}(#1)}
\providecommand{\Deduce}{\mathrm{D}}
\providecommand{\Dual}{\mathrm{d}}
\providecommand{\Ax}{\mathrm{Ax}}
\providecommand{\pSat}[3]{#1 \models #2\,[#3]}
\providecommand{\pAssign}{\sigma}
\providecommand{\bforall}{\forall}
\providecommand{\bexists}{\exists}
\providecommand{\liff}{\mathbin{\leftrightarrow}}
\providecommand{\FV}{\operatorname{FV}}
\providecommand{\UnaryInf}{\ensuremath{\mid}}
\providecommand{\BinaryInf}{\ensuremath{\mid}}
\providecommand{\depth}{\operatorname{depth}}
\providecommand{\oliflabeldef}[3]{#1}
\providecommand{\TMblank}{\sqcup}
\providecommand{\TMstroke}{\mid}
\providecommand{\TMstay}{\mathrm{S}}
\providecommand{\TMright}{\mathrm{R}}
\providecommand{\TMleft}{\mathrm{L}}
\providecommand{\TMendtape}{\triangleright}
\providecommand{\umin}{\mu}
\providecommand{\Sum}{\sum}
\providecommand{\citep}[1]{\cite{#1}}
\providecommand{\citet}[1]{\cite{#1}}
\providecommand{\citeauthor}[1]{\cite{#1}}
\providecommand{\citeyear}[1]{\cite{#1}}

% ===================== OLP compatibility layer (3) =====================
% --- third compatibility block: remaining notation seen in the frozen OLP units ---
\providecommand{\iso}{\cong}
\providecommand{\nszero}{\mathrm{0}}
\providecommand{\nssucc}{\mathrm{S}}
\providecommand{\nsplus}{\mathbin{+}}
\providecommand{\nssless}{<}
\providecommand{\nsless}{<}
\providecommand{\PL}{\mathrm{PL}}
\providecommand{\Nec}{\mathord{\Box}}
\providecommand{\RK}{\mathrm{RK}}
\providecommand{\ST}{\mathrm{ST}}
\providecommand{\SSubst}[3]{#1[#2/#3]}
\providecommand{\subst}[3]{#1[#2/#3]}
\providecommand{\tagfalse}{\mathrm{F}}
\providecommand{\cutr}{\mathrm{Cut}}
\providecommand{\cardneq}{\neq}
\providecommand{\issucc}[1]{\operatorname{succ}(#1)}
\providecommand{\tup}[1]{\langle #1 \rangle}
\providecommand{\Var}{\mathrm{Var}}
\providecommand{\Bew}{\operatorname{Bew}}
\providecommand{\Con}{\operatorname{Con}}
\providecommand{\Cons}{\operatorname{Cons}}
\providecommand{\pDeduce}{\mathrm{D}}
\providecommand{\psat}{\models}
\providecommand{\den}[1]{\llbracket #1 \rrbracket}
\providecommand{\Sem}[2]{\llbracket #1 \rrbracket_{#2}}
\providecommand{\DomainOf}[1]{\lvert #1 \rvert}
\providecommand{\Form}{\mathrm{Form}}
\providecommand{\Terms}{\mathrm{Term}}
\providecommand{\Atom}{\mathrm{Atom}}
\providecommand{\FVars}[1]{\operatorname{FV}(#1)}
\providecommand{\fv}[1]{\operatorname{fv}(#1)}
\providecommand{\tuplep}[1]{\langle #1 \rangle}
\providecommand{\restrict}[2]{#1 \upharpoonright #2}
\providecommand{\sto}{\rightarrow}
\providecommand{\FmlaSet}[1]{\mathrm{Fmla}_{#1}}
\providecommand{\Thm}{\mathrm{Thm}}
\providecommand{\Eval}{\operatorname{val}}
\providecommand{\blank}{\sqcup}
\providecommand{\blankTape}{\sqcup}
\providecommand{\Consis}{\operatorname{Con}}
\providecommand{\drs}{\mathrm{d}}
\providecommand{\VRule}{\mid}
\providecommand{\Bet}{\mathrm{B}}
\providecommand{\Seg}[1]{\mathrm{Seg}_{#1}}
\providecommand{\Enums}{\mathrm{En}}
\providecommand{\Z}{\mathbb{Z}}
\providecommand{\N}{\mathbb{N}}
\providecommand{\abs}[1]{\lvert #1 \rvert}
\providecommand{\floor}[1]{\lfloor #1 \rfloor}
\providecommand{\ceil}[1]{\lceil #1 \rceil}

% ===================== OLP compatibility layer (4) =====================
% --- fourth compatibility block: harvested from the per-candidate compile logs ---
\providecommand{\beredpar}{\mathrel{\rightarrow_\beta}}
\makeatletter
\def\bered{\@ifnextchar[{\bered@i}{\mathrel{\rightarrow_\beta}}}
\def\bered@i[#1]{\mathrel{\rightarrow_\beta}}
\def\bred{\@ifnextchar[{\bred@i}{\mathrel{\rightarrow_\beta}}}
\def\bred@i[#1]{\mathrel{\rightarrow_\beta}}
\makeatother
\providecommand{\beredone}{\mathrel{=_\beta}}
\providecommand{\bredone}{\mathrel{=_\beta}}
\providecommand{\xredone}{\mathrel{=_\beta}}
\providecommand{\aconvone}{\mathrel{=_\alpha}}
\providecommand{\cardless}{<}
\providecommand{\tagrefs}[1]{#1}
\providecommand{\ordplus}{\mathbin{+}}
\providecommand{\funrestrictionto}[2]{#1 \upharpoonright #2}
\providecommand{\fundefined}{\mathord{\uparrow}}
\providecommand{\Trm}{\mathrm{Trm}}
\providecommand{\OProv}{\operatorname{Prov}}
\providecommand{\OPrf}{\operatorname{Prf}}
\providecommand{\aeq}{\mathrel{=}}
\providecommand{\mClass}[1]{[#1]}
\providecommand{\equivclass}[1]{[#1]}
\providecommand{\pheight}{\operatorname{ht}}
\providecommand{\phy}{\operatorname{ht}}
\providecommand{\Cut}{\mathrm{Cut}}
\providecommand{\QuantRank}{\operatorname{qr}}
\providecommand{\elemequiv}{\mathrel{\equiv}}
\providecommand{\PIso}{\mathrm{PIso}}
\providecommand{\Theory}{\mathrm{Th}}
\providecommand{\supstrict}{\subsetneq}
\providecommand{\trcl}{\operatorname{trcl}}
\providecommand{\cref}[1]{\ref{#1}}
\providecommand{\Discharge}{\mathrm{Dis}}
\providecommand{\Weakening}{\mathrm{Wk}}
\providecommand{\ZFminus}{\mathrm{ZF}^{-}}
\makeatletter
\renewcommand{\Dual}{\@ifnextchar[{\Dual@i}{\ifmmode\mathrm{d}\else$\mathrm{d}$\fi}}
\def\Dual@i[#1]{\ifmmode\mathrm{d}[#1]\else$\mathrm{d}[#1]$\fi}
\makeatother
\providecommand{\RightR}{\mathrm{R}}
\providecommand{\LeftR}{\mathrm{L}}
\providecommand{\Intro}{\mathrm{I}}
\providecommand{\Axiom}{\mathrm{Ax}}
\providecommand{\formula}[1]{#1}
\providecommand{\stared}{\ast}
\providecommand{\oset}[2]{\{\, #1 \mid #2 \,\}}
\providecommand{\setof}[2]{\{\, #1 \mid #2 \,\}}
\providecommand{\pow}[1]{\mathcal{P}(#1)}
\providecommand{\ordpair}[2]{\langle #1, #2 \rangle}
\providecommand{\ordtriple}[3]{\langle #1, #2, #3 \rangle}
\providecommand{\union}{\cup}
\providecommand{\inter}{\cap}
\providecommand{\seq}[1]{\langle #1 \rangle}
\providecommand{\SuchThat}{\mid}
\providecommand{\where}{\mid}
\providecommand{\TVar}{\mathrm{TVar}}
\providecommand{\Bool}{\mathrm{Bool}}
\providecommand{\Val}{\operatorname{Val}}
\providecommand{\TL}[1]{\mathrm{TL}}
\providecommand{\xstared}{\ast}

% ===================== OLP compatibility layer (5) =====================
% --- fifth compatibility block: rescues the quarantined OLP candidates ---
\providecommand{\ident}{\operatorname{id}}
\providecommand{\setrank}{\operatorname{rank}}
\providecommand{\cardfont}{\mathfrak{c}}
\providecommand{\Complement}{\overline}
\providecommand{\QR}{\mathrm{QR}}
\providecommand{\becd}{\ensuremath{\ulcorner}}
\providecommand{\disjointsum}{\mathbin{\oplus}}
\providecommand{\cardeq}[2]{#1 = #2}
\providecommand{\xred}{\mathrel{\rightarrow}}
% --- sixth compatibility block: names reported by real compile logs ---
\providecommand{\ORefut}{\operatorname{Refut}}
\providecommand{\ORProv}{\operatorname{Prov}}
\providecommand{\cutrank}{\operatorname{cr}}
\providecommand{\Refut}{\operatorname{Refut}}
\providecommand{\DischargeRule}{\mathrm{Dis}}
\providecommand{\Contraction}{\mathrm{C}}
\providecommand{\bottomAlignProof}{\relax}
\providecommand{\FalseCl}{\bot}
\providecommand{\red}{\mathrel{\rightarrow}}
\providecommand{\varAssign}{\sigma}
\providecommand{\defiff}{\mathrel{\leftrightarrow}}
\providecommand{\FalseInt}{\bot}
\providecommand{\fn}{\operatorname{fn}}
\providecommand{\dotted}{\cdot}

% --- tag-macro semantic fix (2026-10-13, final): \iftag CHOOSES a branch and the OLP proof body is
% always the SECOND group (\iftag{tag}{body}, {tag}{body}{c}, {tag}{body}{c}{d}); a definition that
% prints the tag literal drops the proof body. Bodies contain blank lines, hence \long.
\makeatletter
\providecommand{\iftag}{}
\long\def\iftag#1#2{\@ifnextchar\bgroup{\iftag@three{#2}}{\iftag@stop{#2}}}
\long\def\iftag@three#1#2{\iftag@stop{#1}\@ifnextchar\bgroup{\iftag@eat}{\relax}}
\long\def\iftag@eat#1{}
\long\def\iftag@stop#1{#1}
\long\def\oliflabeldef#1#2{\@ifnextchar\bgroup{\@gobble}{\iftag@stop{#2}}}
\providecommand{\tagrefs}[1]{}
\providecommand{\settag}[1]{}
\makeatother
\begin{document}

% ==== OLP file: content/normal-modal-logic/frame-definability/definability.tex ====
% Part: normal-modal-logic
% Chapter: frame-definability
% Section: definability



\olfileid{nml}{frd}{def}

\olsection{Frame Definability}

Even though the converse implications of \olref[acc]{thm:soundschemas}
fail, they hold if we replace ``model'' by ``frame'': for the
properties considered in \olref[acc]{thm:soundschemas}, it \emph{is}
true that if \olup{formula} is valid in a \emph{frame} then the
accessibility relation of that frame has the corresponding
property. So, the \olup{formula}s considered \emph{define} the classes of
frames that have the corresponding property.

\begin{defn}
  If $\mClass{F}$ is a class of frames, we say $A$
  \emph{defines}~$\mClass{F}$ iff $\mModel{F} \Entails A$ for all and only
  frames~$\mModel{F} \in \mClass{F}$.
\end{defn}

We now proceed to establish the full definability results for
frames.

\begin{thm}\ollabel{thm:fullCorrespondence}
If the \olup{formula} on the right side of \olref[acc]{tab:five} is valid in a
frame~$\mModel{F}$, then $\mModel{F}$ has the property on the left
side.
\end{thm}

\begin{proof}
  \begin{enumerate}
  \item Suppose \Ax{D} is valid in $\mModel{F} = \tuple{W, R}$, i.e.,
    $\mModel{F} \Entails \Box p \lif \Diamond p$. Let $\mModel{M} =
    \tuple{W, R, V}$ be a model based on~$\mModel{F}$, and $w \in
    W$. We have to show that there is a $v$ such that~$Rwv$.  Suppose
    not: then both $\mSat{M}{\Box A}$ and $\mSat/{M}{\Diamond A}[w]$
    for any~$A$, including~$p$. But then $\mSat/{M}{\Box p \lif
      \Diamond p}[w]$, contradicting the assumption that $\mModel{F}
    \Entails \Box p \lif \Diamond p$.
  \item Suppose \Ax{T} is valid in $\mModel{F}$, i.e., $\mModel{F} \Entails \Box
    p \lif p$. Let $w \in W$ be an arbitrary world; we need to show
    $Rww$. Let $u \in V(p)$ if and only if $Rwu$ (when $q$ is other
    than $p$, $V(q)$ is arbitrary, say $V(q) = \emptyset)$. Let
    $\mModel{M} = \tuple{W, R, V}$. By construction, for all $u$ such
    that $Rwu$: $\mSat{M}{p}[u]$, and hence $\mSat{M}{\Box p}[w]$. But
    by hypothesis $\Box p \lif p$ is true at~$w$, so that
    $\mSat{M}{p}[w]$, but by definition of $V$ this is possible only
    if~$Rww$.
  \item We prove the contrapositive: Suppose $\mModel{F}$ is not
    symmetric, we show that \Ax{B}, i.e., $p \lif \Box\Diamond p$ is
    not valid in $\mModel{F}= \tuple{W, R}$. If $\mModel{F}$ is not
    symmetric, there are $u$, $v \in W$ such that $Ruv$ but not
    $Rvu$. Define $V$ such that $w \in V(p)$ if and only if not $Rvw$
    (and $V$ is arbitrary otherwise). Let $\mModel{M} =\tuple{W, R,
      V}$.  Now, by definition of $V$, $\mSat{M}{p}[w]$ for all $w$
    such that not $Rvw$, in particular, $\mSat{M}{p}[u]$ since
    not~$Rvu$. Also, since $Rvw$ iff $w \notin V(p)$, there is no~$w$
    such that $Rvw$ and $\mSat{M}{p}[w]$, and hence
    $\mSat/{M}{\Diamond p}[v]$. Since~$Ruv$, also
    $\mSat/{M}{\Box\Diamond p}[u]$. It follows that $\mSat/{M}{p \lif
      \Box\Diamond p}[u]$, and so $\Ax{B}$ is not valid
    in~$\mModel{F}$.
  \item Suppose \Ax{4} is valid in $\mModel{F} = \tuple{W,R}$, i.e.,
    $\mModel{F} \Entails \Box p \lif \Box\Box p$, and let $u$, $v$, $w \in W$ be
    arbitrary worlds such that $Ruv$ and $Rvw$; we need to show that
    $Ruw$. Define $V$ such that $z \in V(p)$ if and only if $Ruz$ (and
    $V$ is arbitrary otherwise). Let $\mModel{M} =\tuple{W, R, V}$. By
    definition of $V$, $\mSat{M}{p}[z]$ for all $z$ such that $Ruz$,
    and hence $\mSat{M}{\Box p}[u]$. But by hypothesis \Ax{4}, $\Box p
    \lif \Box \Box p$, is true at $u$, so that $\mSat{M}{\Box \Box
      p}[u]$. Since $Ruv$ and $Rvw$, we have $\mSat{M}{p}[w]$, but by
    definition of~$V$ this is possible only if $Ruw$, as desired.
  \item We proceed contrapositively, assuming that the frame
    $\mModel{F} = \tuple{W, R}$ is not euclidean, and show that it
    falsifies~\Ax{5}, i.e., $\mModel{F} \Entails/ \Diamond p \lif
      \Box\Diamond p$. Suppose there are worlds $u$, $v$, $w \in W$
    such that $Rwu$ and $Rwv$ but not~$Ruv$. Define $V$ such that for
    all worlds $z$, $z \in V(p)$ if and only if it is \emph{not} the
    case that~$Ruz$. Let $\mModel{M} =\tuple{W, R, V}$. Then by
    hypothesis $\mSat{M}{p}[v]$ and since $Rwv$ also
    $\mSat{M}{\Diamond p}[w]$. However, there is no world $y$ such
    that $Ruy$ and $\mSat{M}{p}[y]$ so $\mSat/{M}{\Diamond
      p}[u]$. Since $Rwu$, it follows that $\mSat/{M}{\Box\Diamond
      p}[w]$, so that \Ax{5}, $\Diamond p \lif \Box\Diamond p$, fails
    at~$w$.
  \end{enumerate}
\end{proof}

You'll notice a difference between the proof for \Ax{D} and the other
cases: no mention was made of the valuation~$V$. In effect, we proved
that if $\mSat{M}{\Ax{D}}$ then $\mModel{M}$ is serial. So \Ax{D}
defines the class of serial \emph{models}, not just frames.

\begin{cor}\ollabel{prop:D-serial}
  Any model where \Ax{D} is true is serial.
\end{cor}

\begin{cor}
Each \olup{formula} on the right side of \olref[acc]{tab:five} defines
the class of frames which have the property on the left side.
\end{cor}

\begin{proof}
  In \olref[acc]{thm:soundschemas}, we proved that if a model has the
  property on the left, the \olup{formula} on the right is true in
  it. Thus, if a frame~$\mModel{F}$ has the property on the left, the
  \olup{formula} on the right is valid in~$\mModel{F}$. In
  \olref{thm:fullCorrespondence}, we proved the converse implications:
  if \olup{formula} on the right is valid in~$\mModel{F}$, $\mModel{F}$
  has the property on the left.
\end{proof}

\begin{prob}
Show that if the \olup{formula} on the right side of
\olref[nml][frd][acc]{tab:anotherfive} is valid in a frame~$\mModel{F}$, then
$\mModel{F}$ has the property on the left side. To do this, consider a
frame that does \emph{not} satisfy the property on the left, and
define a suitable~$V$ such that the \olup{formula} on the right is false
at some world.
\end{prob}

\olref{thm:fullCorrespondence} also shows that the properties can be
combined: for instance if both \Ax{B} and \Ax{4} are valid
in~$\mModel{F}$ then the frame is both symmetric and transitive,
etc. Many important modal logics are characterized as the set of
\olup{formula}s valid in all frames that combine some frame properties,
and so we can characterize them as the set of \olup{formula}s valid in
all frames in which the corresponding defining \olup{formula}s are valid.
For instance, the classical system \Log{S4} is the set of all
\olup{formula}s valid in all reflexive and transitive frames, i.e., in
all those where both \Ax{T} and~\Ax{4} are valid.  \Log{S5} is the set
of all formulas valid in all reflexive, symmetric, and euclidean
frames, i.e., all those where all of \Ax{T}, \Ax{B}, and \Ax{5} are
valid.

Logical relationships between properties of $R$ in general correspond
to relationships between the corresponding defining \olup{formula}s. For
instance, every reflexive relation is serial; hence, whenever \Ax{T}
is valid in a frame, so is~\Ax{D}. (Note that this relationship is
\emph{not} that of entailment. It is not the case that whenever
$\mSat{M}{\Ax{T}}[w]$ then $\mSat{M}{\Ax{D}}[w]$.) We record some such
relationships.

\begin{prop}\ollabel{prop:relation-facts}
  Let $R$ be a binary relation on a set $W$; then:
  \begin{enumerate}
  \item If $R$ is reflexive, then it is serial.
  \item If $R$ is symmetric, then it is transitive if and only if it
    is euclidean.
  \item If $R$ is symmetric or euclidean then it is weakly directed
    (it has the ``diamond property'').
  \item If $R$ is euclidean then it is weakly connected.
  \item If $R$ is functional then it is serial.  
  \end{enumerate}
\end{prop}

\begin{prob}
  Prove \olref[nml][frd][def]{prop:relation-facts}.
\end{prob}


% ==== OLP file: content/normal-modal-logic/frame-definability/equivalence-S5.tex ====
% Part: normal-modal-logic
% Chapter: frame-correspondence
% Section: equivalence-S5



\olfileid{nml}{frd}{es5}

\olsection{Equivalence Relations and \Log{S5}}

The modal logic \Log{S5} is characterized as the set of \olup{formula}s
valid on all universal frames, i.e., every world is accessible from
every world, including itself. In such a scenario, $\Box$ corresponds
to necessity and $\Diamond$ to possibility: $\Box A$ is true if $A$
is true at \emph{every} world, and $\Diamond A$ is true if $A$ is
true at \emph{some} world. It turns out that \Log{S5} can also be
characterized as the \olup{formula}s valid on all reflexive, symmetric,
and transitive frames, i.e., on all \emph{equivalence relations}.

\begin{defn}
  A binary relation $R$ on $W$ is an \emph{equivalence relation} if
  and only if it is reflexive, symmetric and transitive.  A relation
  $R$ on $W$ is \emph{universal} if and only if $Ruv$ for all $u,v \in
  W$.
\end{defn}

Since \Ax{T}, \Ax{B}, and \Ax{4} characterize the reflexive,
symmetric, and transitive frames, the frames where the accessibility
relation is an equivalence relation are exactly those in which all
three \olup{formula}s are valid. It turns out that the equivalence
relations can also be characterized by other combinations of
\olup{formula}s, since the conditions with which we've defined
equivalence relations are equivalent to combinations of other familiar
conditions on~$R$.

\begin{prop}\ollabel{prop:equivalences}
  The following are equivalent:
  \begin{enumerate}
  \item $R$ is an equivalence relation;
  \item $R$ is reflexive and euclidean;
  \item $R$ is serial, symmetric, and euclidean;
  \item $R$ is serial, symmetric, and transitive.
  \end{enumerate}
\end{prop}

\begin{proof}
  Exercise.
\end{proof}

\begin{prob}
  Prove \olref[nml][frd][es5]{prop:equivalences} by showing:
  \begin{enumerate}
  \item If $R$ is symmetric and transitive, it is euclidean.
  \item If $R$ is reflexive, it is serial.
  \item If $R$ is reflexive and euclidean, it is symmetric.
  \item If $R$ is symmetric and euclidean, it is transitive.
  \item If $R$ is serial, symmetric, and transitive, it is reflexive.
  \end{enumerate}
  Explain why this suffices for the proof that the conditions are
  equivalent.
\end{prob}

\olref{prop:equivalences} is the semantic counterpart to
\olref[prf][prs]{prop:S5}, in that it gives an equivalent characterization of
the modal logic of frames over which $R$ is an equivalence relation (the logic
traditionally referred to as \Log{S5}).

What is the relationship between universal and equivalence relations?
Although every universal relation is an equivalence relation, clearly
not every equivalence relation is universal. However, the \olup{formula}s
valid on all universal relations are exactly the same as those valid
on all equivalence relations.

\begin{prop}
  Let $R$ be an equivalence relation, and for each $w \in W$ define
  the \emph{equivalence class} of $w$ as the set $[w] = \{w'\in W :
  Rww'\}$. Then:
  \begin{enumerate}
  \item $w \in [w]$;
  \item $R$ is universal on each equivalence class $[w]$;
  \item The collection of equivalence classes partitions $W$ into mutually
    exclusive and jointly exhaustive subsets.
  \end{enumerate}
\end{prop}

\begin{prop}\ollabel{prop:S5=univ}
  !!^a{formula}~$A$ is valid in all frames $\mModel{F} =\tuple{W,R}$
  where $R$ is an equivalence relation, if and only if it is valid in
  all frames $\mModel{F} =\tuple{W,R}$ where $R$ is universal. Hence,
  the logic of universal frames is just \Log{S5}.
\end{prop}

\begin{proof}
  It's immediate to verify that a universal relation~$R$ on~$W$ is an
  equivalence. Hence, if $A$ is valid in all frames where $R$ is an
  equivalence it is valid in all universal frames. For the other
  direction, we argue contrapositively: suppose $B$ is \olup{formula}
  that fails at a world $w$ in a model $\mModel{M} = \tuple{W, R, V}$
  based on a frame $\tuple{W,R}$, where $R$ is an equivalence on
  $W$. So $\mSat/{M}{B}[w]$. Define a model $\mModel{M}' = \tuple{W',
    R', V'}$ as follows:
  \begin{enumerate}
  \item $W' = [w]$;
  \item $R'$ is universal on $W'$;
  \item $V'(p) = V(p) \cap W'$. 
  \end{enumerate}
  (So the set $W'$ of worlds in $\mModel{M}'$ is represented by the
  shaded area in \olref{fig:partition}.)  It is easy to see
  that $R$ and $R'$ agree on $W'$. Then one can show by induction on
  \olup{formula}s that for all $w' \in W'$: $\mSat{M'}{A}[w']$ if and
  only if $\mSat{M}{A}[w']$ for each $A$ (this makes sense since $W'
  \subseteq W$). In particular, $\mSat/{M'}{B}[w]$, and $B$ fails in
  a model based on a universal frame.
\end{proof}

\begin{figure}[t]
  \centering
  \begin{tikzpicture}[node distance=2cm, auto, thick]
    \clip [rounded corners] (0,0) -- (6,0) -- (6,4) -- (0,4) --  cycle;
    \begin{scope}
      \clip (0,0) -- (2,0)  .. controls (1.5,1.5) and (4.5,1.5)
      .. (4,4) -- (0,4) -- (0,0) -- cycle;
      \filldraw[fill=gray!40] (0,4) -- (0,2) .. controls (1.5,1.5)
      and (4.5,1.5) .. (4.5,0) -- (6,0) -- (6,4) -- (0,4) --  cycle;
    \end{scope}
    \draw [name path=line1] (2,0) .. controls (1.5,1.5) and (4.5,1.5) .. (4,4);
    \draw [name path=line2] (0,2)  .. controls (1.5,1.5) and (4.5,1.5) .. (4.5,0);
    \path [name intersections={of=line1 and line2}];
    \draw [rounded corners, use as bounding box, very thick] (0,0)
    -- (6,0) -- (6,4) -- (0,4) --  cycle; 
    \node at (2,3) {$[w]$} ;
    \node at (1,1) {$[u]$} ;
    \node at (3,0.75) {$[v]$} ;
    \node at (4.75,2) {$[z]$} ;
  \end{tikzpicture}
  \caption{A partition of $W$ in equivalence classes.}
  \ollabel{fig:partition}
\end{figure}



% ==== OLP file: content/normal-modal-logic/frame-definability/first-order-definability.tex ====
% Part: normal-modal-logic
% Chapter: frame-correspondence
% Section: first-order-definability



\olfileid{nml}{frd}{fol}

\olsection{First-order Definability}

We've seen that a number of properties of accessibility relations of
frames can be defined by modal \olup{formula}s. For instance, symmetry of
frames can be defined by the \olup{formula}~\Ax{B}, $p \lif \Box\Diamond
p$. The conditions we've encountered so far can all be expressed by
first-order \olup{formula}s in \olup{language} involving a single two-place
\olup{predicate}. For instance, symmetry is defined by
$\lforall[x][\lforall[y][(\Atom{Q}{x,y} \lif \Atom{Q}{y, x})]]$ in the
sense that a first-order \olup{structure}~$\Struct{M}$ with $\Domain{M} =
W$ and $\Assign{Q}{M} = R$ satisfies the preceding formula iff $R$ is
symmetric. This suggests the following definition:

\begin{defn}
  A class~$\mClass{F}$ of frames is \emph{first-order definable} if there is
  \olup{sentence}~$A$ in the first-order language with a single
  two-place \olup{predicate}~$Q$ such that $\mModel{F} =
  \tuple{W, R} \in \mClass{F}$ iff $\Sat{M}{A}$ in the first-order
  \olup{structure}~$\Struct{M}$ with $\Domain{M} = W$ and $\Assign{Q}{M}
  = R$.
\end{defn}

It turns out that the properties and modal \olup{formula}s that define
them considered so far are exceptional. Not every \olup{formula} defines
a first-order definable class of frames, and not every first-order
definable class of frames is definable by a modal formula.

A counterexample to the first is given by the L\"ob formula:
\begin{equation}
\Box(\Box p \lif p) \lif \Box p. \tag{\Ax{W}}
\end{equation}
\Ax{W} defines the class of transitive and converse well-founded
frames. A relation is well-founded if there is no infinite sequence
$w_1$, $w_2$, \dots{} such that $Rw_2w_1$, $Rw_3w_2$, \dots. For
instance, the relation $<$ on~$\Nat$ is well-founded, whereas the
relation $<$ on~$\Int$ is not. A relation is converse well-founded
iff its converse is well-founded. So converse well-founded relations
are those where there is no infinite sequence $w_1$, $w_2$, \dots{}
such that $Rw_1w_2$, $Rw_2w_3$, \dots. 

There is, however, no first-order \olup{formula} defining transitive
converse well-founded relations. For suppose~$\Sat{M}{F}$ iff $R =
\Assign{Q}{M}$ is transitive converse well-founded. Let $A_n$ be the
\olup{formula}
\[
(\Atom{Q}{a_1,a_2} \land \dots \land
    \Atom{Q}{a_{n-1},a_{n}})
\]
Now consider the set of \olup{formula}s
\[
\Gamma = \{F, A_1, A_2, \dots\}.
\]
Every finite subset of $\Gamma$ is satisfiable: Let $k$ be largest
such that $A_k$ is in the subset, $\Domain{M_k} = \{1, \dots, k\}$,
$\Assign{a_i}{M_k} = i$, and $\Assign{Q}{M_k} = <$. Since $<$ on $\{1,
\dots, k\}$ is transitive and converse well-founded,
$\Sat{M_k}{F}$. $\Sat{M_k}{A_i}$ by construction, for all $i \le k$.
By the Compactness Theorem for first-order logic, $\Gamma$ is
satisfiable in some \olup{structure}~$\Struct{M}$. By hypothesis, since
$\Sat{M}{F}$, the relation $\Assign{Q}{M}$ is converse
well-founded. But clearly, $\Assign{a_1}{M}$, $\Assign{a_2}{M}$,
\dots{} would form an infinite sequence of the kind ruled out by
converse well-foundedness.

A counterexample to the second claim is given by the property of
universality: for every $u$ and $v$, $Ruv$. Universal frames are
first-order definable by the formula
$\lforall[x][\lforall[y][\Atom{Q}{x,y}]]$.  However, no modal formula
is valid in all and only the universal frames. This is a consequence
of a result that is independently interesting: the formulas valid in
universal frames are exactly the same as those valid in reflexive,
symmetric, and transitive frames. There are reflexive, symmetric, and
transitive frames that are not universal, hence every \olup{formula}
valid in all universal frames is also valid in some non-universal
frames.


% ==== OLP file: content/normal-modal-logic/frame-definability/frame-definability.tex ====
% Part: normal-modal-logic
% Chapter: frame-definability



\olchapter{nml}{frd}{Frame Definability}

\olimport{introduction}
\olimport{properties-accessibility}
\olimport{frames}
\olimport{definability}
\olimport{first-order-definability}
\olimport{equivalence-S5}
\olimport{second-order-definability}

\OLEndChapterHook



% ==== OLP file: content/normal-modal-logic/frame-definability/frames.tex ====
% Part: normal-modal-logic
% Chapter: frame-correspondence
% Section: frames



\olfileid{nml}{frd}{fra}

\olsection{Frames}

\begin{defn}
  A \emph{frame} is a pair $\mModel{F} = \tuple{W,R}$ where $W$ is a
  non-empty set of worlds and $R$ a binary relation on~$W$. A model
  $\mModel{M}$ is \emph{based on} a frame $\mModel{F} = \tuple{W,R}$
  if and only if $\mModel{M} = \tuple{W, R, V}$ for some valuation $V$.
\end{defn}

\begin{defn}
  If $\mModel{F}$ is a frame, we say that $A$ is \emph{valid in
    $\mModel{F}$,} $\mModel{F} \Entails A$, if $\mSat{M}{A}$ for
  every model~$\mModel{M}$ based on~$\mModel{F}$.
  
  If $\mClass{F}$ is a class of frames, we say $A$ is \emph{valid in
    $\mClass{F}$,} $\mClass{F} \Entails A$, iff $\mModel{F} \Entails
  A$ for every frame~$\mModel{F} \in \mClass{F}$.
\end{defn}

The reason frames are interesting is that correspondence between
schemas and properties of the accessibility relation~$R$ is at the
level of frames, \emph{not of models}. For instance, although \Ax{T}
is true in all reflexive models, not every model in which \Ax{T} is
true is reflexive. However, it \emph{is} true that not only is \Ax{T}
\emph{valid} on all reflexive \emph{frames}, also every frame in
which \Ax{T} is valid is reflexive.

\begin{rem}
Validity in a class of frames is a special case of the notion of
validity in a class of models: $\mClass{F} \Entails A$ iff
$\mClass{C} \Entails A$ where $\mClass{C}$ is the class of all models
based on a frame in~$\mClass{F}$.

Obviously, if \olup{formula} or a schema is valid, i.e., valid with
respect to the class of \emph{all} models, it is also valid with
respect to any class~$\mClass{F}$ of frames.
\end{rem}


% ==== OLP file: content/normal-modal-logic/frame-definability/introduction.tex ====
% Part: normal-modal-logic
% Chapter: frame-correspondence
% Section: introduction



\olfileid{nml}{frd}{int}

\olsection{Introduction}

One question that interests modal logicians is the relationship
between the accessibility relation and the truth of certain
\olup{formula}s in models with that accessibility relation. For instance,
suppose the accessibility relation is reflexive, i.e., for every $w
\in W$, $Rww$. In other words, every world is accessible from
itself. That means that when $\Box A$ is true at a world~$w$, $w$
itself is among the accessible worlds at which~$A$ must therefore be
true. So, if the accessibility relation~$R$ of~$\mModel{M}$ is
reflexive, then whatever world $w$ and formula $A$ we take, $\Box A
\lif A$ will be true there (in other words, the schema $\Box p \lif
p$ and all its substitution instances are true in~$\mModel{M}$).

The converse, however, is false. It's not the case, e.g., that if
$\Box p \lif p$ is true in~$\mModel{M}$, then $R$ is reflexive. For we
can easily find a non-reflexive model~$\mModel{M}$ where $\Box p \lif
p$ is true at all worlds: take the model with a single world~$w$, not
accessible from itself, but with $w \in V(p)$. By picking the truth
value of $p$ suitably, we can make $\Box A \lif A$ true in a model
that is not reflexive.

The solution is to remove the variable assignment~$V$ from the
equation. If we require that $\Box p \lif p$ is true at all worlds
in~$\mModel{M}$, regardless of which worlds are in~$V(p)$, then it is
necessary that $R$ is reflexive. For in any non-reflexive model, there
will be at least one world~$w$ such that not $Rww$. If we set $V(p) =
W \setminus \{w\}$, then $p$ will be true at all worlds other
than~$w$, and so at all worlds accessible from~$w$ (since $w$ is
guaranteed not to be accessible from~$w$, and $w$ is the only world
where~$p$ is false). On the other hand, $p$ is false at $w$, so $\Box
p \lif p$ is false at~$w$.

This suggests that we should introduce a notation for model structures
without a valuation: we call these \emph{frames}. A frame $\mModel{F}$
is simply a pair $\tuple{W, R}$ consisting of a set of worlds with an
accessibility relation. Every model $\tuple{W, R, V}$ is then, as we
say, \emph{based on} the frame $\tuple{W, R}$. Conversely, a frame
determines the class of models based on it; and a class of frames
determines the class of models which are based on any frame in the
class. And we can define $\mModel{F} \Entails A$, the notion of \olup{formula}
being \emph{valid} in a frame as: $\mSat{M}{A}$ for all $\mModel{M}$
based on~$\mModel{F}$.

With this notation, we can establish correspondence relations between
\olup{formula}s and classes of frames: e.g., $\mModel{F} \Entails \Box p \lif p$
if, and only if, $\mModel{F}$ is reflexive.


% ==== OLP file: content/normal-modal-logic/frame-definability/properties-accessibility.tex ====
% Part: normal-modal-logic
% Chapter: frame-correspondence
% Section: properties-accessibility



\olfileid{nml}{frd}{acc}

\olsection{Properties of Accessibility Relations}

Many modal \olup{formula}s turn out to be characteristic of simple, and
even familiar, properties of the accessibility relation. In one
direction, that means that any model that has a given property makes a
corresponding \olup{formula} (and all its substitution instances)
true. We begin with five classical examples of kinds of accessibility
relations and the formulas the truth of which they guarantee.

\begin{thm}\ollabel{thm:soundschemas}
  Let $\mModel{M} = \tuple{W, R, V}$ be a model.  If $R$ has the
  property on the left side of \olref{tab:five}, every instance of the
  \olup{formula} on the right side is true in~$\mModel{M}$.
\end{thm}

\begin{table}[t]
    \begin{tabular}{| l || l |}
      \hline
      {\emph{If $R$ is \dots}} & {\emph{then \dots is true in~$\mModel{M}$:}} \\
      \hline \hline
      \emph{serial}: $\forall u \exists v Ruv$ & \hbox to.43\textwidth
           {$\Box p \lif \Diamond p$ \hfill (\Ax{D})} \\
      \hline
      \emph{reflexive}: $\forall w Rww$  
      & \hbox to.43\textwidth
          {$\Box p \lif p$ \hfill (\Ax{T})} \\
      \hline
      \emph{symmetric}: &  \hbox to .43\textwidth
           {$p \lif \Box\Diamond p$ \hfill (\Ax{B})} \\
           $\forall u\forall v(Ruv \lif Rvu)$ & \\
      \hline
      \emph{transitive}: & \hbox to.43\textwidth
           {$\Box p \lif \Box \Box p$ \hfill (\Ax{4})} \\
      $\forall u \forall v \forall w ((Ruv \land Rvw) \lif Ruw)$ & \\
      \hline 
      \emph{euclidean}: & \hbox to .43\textwidth
        {$\Diamond p \lif \Box\Diamond p$ \hfill (\Ax{5})} \\
      $\forall w \forall u \forall v ((Rwu \land Rwv) \lif Ruv)$ &\\
      \hline
    \end{tabular}
    \caption{Five correspondence facts.}
    \ollabel{tab:five}
  \end{table} 


\begin{proof}
  Here is the case for \Ax{B}: to show that the schema is true in a
  model we need to show that all of its instances are true at all worlds
  in the model. So let $A \lif \Box\Diamond A$ be a given instance
  of \Ax{B}, and let $w \in W$ be an arbitrary world. Suppose the
  antecedent $A$ is true at $w$, in order to show that $\Box \Diamond
  A$ is true at $w$. So we need to show that $\Diamond A$ is true at
  all $w'$ accessible from $w$. Now, for any $w'$ such that $Rww'$ we
  have, using the hypothesis of symmetry, that also $Rw'w$ (see
  \olref{fig:Bsymm}). Since $\mSat{M}{A}[w]$, we have
  $\mSat{M}{\Diamond A}[w']$. Since $w'$ was an arbitrary world such
  that $Rww'$, we have $\mSat{M}{\Box\Diamond A}[w]$.

  We leave the other cases as exercises.
\end{proof}

\begin{prob}
  Complete the proof of \olref[nml][frd][acc]{thm:soundschemas}.
\end{prob}

\begin{figure}
  \begin{center}
    \begin{tikzpicture}[modal]
      \node[world] (w1) [label={below:$\mSat{{}}{\formula{A}}$\\
          $\mSat{{}}{\Box\Diamond \formula{A}}$}] {$w$}; 
      \node[world] (w2) [label=below:{$\mSat{{}}{\Diamond \formula{A}}$},
        right=of w1] {$w'$}; 
      \draw[->, bend left] (w1) to (w2); 
      \draw[->, bend left] (w2) to (w1); 
    \end{tikzpicture}
  \end{center}
\caption{The argument from symmetry.}
\ollabel{fig:Bsymm}
\end{figure}

Notice that the converse implications of \olref{thm:soundschemas} do
not hold: it's not true that if a model verifies a schema, then the
accessibility relation of that model has the corresponding property.
In the case of \Ax{T} and reflexive models, it is easy to give an
example of a model in which \Ax{T} itself fails: let $W = \{w\}$ and
$V(p) = \emptyset$. Then $R$ is not reflexive, but $\mSat{M}{\Box
  p}[w]$ and $\mSat/{M}{p}[w]$. But here we have just a single
instance of \Ax{T} that fails in~$\mModel{M}$; other instances, e.g.,
$\Box \lnot p \lif \lnot p$ are true. It is harder to give examples
where \emph{every substitution instance} of~\Ax{T} is true
in~$\mModel{M}$ and $\mModel{M}$ is not reflexive. But there are such
models, too:

\begin{prop}\ollabel{prop:reflexive}
  Let $\mModel{M} = \tuple{W, R, V}$ be a model such that $W = \{u, v
  \}$, where worlds $u$ and $v$ are related by $R$: i.e., both $Ruv$
  and $Rvu$.  Suppose that for all $p$: $u \in V(p) \Leftrightarrow v
  \in V(p)$. Then:
  \begin{enumerate}
  \item For all $A$: $\mSat{M}{A}[u]$ if and only if
    $\mSat{M}{A}[v]$ (use induction on $A$).
  \item Every instance of~\Ax{T} is true in $\mModel{M}$.
  \end{enumerate}
  Since $\mModel{M}$ is not reflexive (it is, in fact,
  \emph{irreflexive}), the converse of \olref{thm:soundschemas}
  fails in the case of \Ax{T} (similar arguments can be given for
  some---though not all---the other schemas mentioned in
  \olref{thm:soundschemas}).
\end{prop}

\begin{prob}
  Prove the claims in \olref[nml][frd][acc]{prop:reflexive}.
\end{prob}

Although we will focus on the five classical \olup{formula}s \Ax{D},
\Ax{T}, \Ax{B}, \Ax{4}, and \Ax{5}, we record in
\olref{tab:anotherfive} a few more properties of accessibility
relations. The accessibility relation~$R$ is partially functional, if
from every world at most one world is accessible. If it is the case
that from every world exactly one world is accessible, we call it
functional. (Thus the functional relations are precisely those that
are both serial and partially functional). They are called
``functional'' because the accessibility relation operates like a
(partial) function. A relation is weakly dense if whenever $Ruv$,
there is a~$w$ ``between'' $u$ and~$v$. So weakly dense relations are
in a sense the opposite of transitive relations: in a transitive
relation, whenever you can reach $v$ from~$u$ by a detour via~$w$, you
can reach $v$ from $u$ directly; in a weakly dense relation, whenever
you can reach $v$ from~$u$ directly, you can also reach it by a detour
via some~$w$. A relation is weakly directed if whenever you can reach
worlds $u$ and~$v$ from some world~$w$, you can reach a single
world~$t$ from both $u$ and~$v$---this is sometimes called the
``diamond property'' or ``confluence.''

\begin{table}[t]
    \begin{tabular}{| p{.50\textwidth} || p{.43\textwidth} |}
      \hline
      {\emph{If $R$ is \dots}} & {\emph{then \dots is true in~$\mModel{M}$:}} \\
      \hline \hline
      \emph{partially functional}: & 
      \multirow{2}{*}{$\Diamond p \lif \Box p$} \\
      $ \forall w \forall u \forall v ((Rwu \land Rwv) \lif u =v )$ & \\
      \hline
      \multirow{2}{*}{\emph{functional}: $\forall w \exists v \forall u(Rwu \liff \eq[u][v]) $} 
      & \multirow{2}{*}{$\Diamond p \liff \Box p$} \\
      & \\
      \hline
      \emph{weakly dense}: &  \multirow{2}{*}{$\Box \Box p \lif
        \Box p$} \\
      $\forall u\forall v(Ruv \lif \exists w(Ruw \land Rwv))$ & \\
      \hline
      \emph{weakly connected}: &
      \multirow{3}{*}{\hbox to.43\textwidth{$
        \begin{array}{@{}l@{}}
          \Box ((p \land \Box p) \lif q) \lor {} \\
          \qquad \Box ((q \land \Box q) \lif p)
        \end{array}$ \hfill (\Ax{L})}
      } \\
      $\forall w \forall u \forall v ((Rwu \land Rwv) \lif {}$ &\\
      \qquad $(Ruv \lor u=v \lor Rvu))$ & \\
      \hline 
      \emph{weakly directed}: & \multirow{3}{*}{\hbox to .43\textwidth
        {$\Diamond\Box p \lif \Box\Diamond p$\hfill (\Ax{G})}} \\
      $\forall w \forall u \forall v ((Rwu \land Rwv) \lif {}$ & \\
      \qquad $\exists t (Rut \land Rvt))$ & \\
      \hline
    \end{tabular}
    \caption{Five more correspondence facts.}
    \ollabel{tab:anotherfive}
  \end{table} 

\begin{prob}
  Let $\mModel{M} = \tuple{W, R, V}$ be a model. Show that if $R$
  satisfies the left-hand properties of
  \olref[nml][frd][acc]{tab:anotherfive}, every instance of the
  corresponding right-hand formula is true in~$\mModel{M}$.
\end{prob}




% ==== OLP file: content/normal-modal-logic/frame-definability/second-order-definability.tex ====
% Part: normal-modal-logic
% Chapter: frame-correspondence
% Section: standard-translation



\olfileid{nml}{frd}{st}

\olsection{Second-order Definability}

Not every frame property definable by modal formulas is first-order
definable. However, if we allow quantification over one-place
predicates (i.e., monadic second-order quantification), we define all
modally definable frame properties. The trick is to exploit a
systematic way in which the conditions under which a modal \olup{formula}
is true at a world are related to first-order \olup{formula}s. This is
the so-called standard translation of modal formulas into first-order
formulas in a language containing not just a two-place
\olup{predicate}~$Q$ for the accessibility relation, but also a one-place
\olup{predicate}~$P_i$ for the \olup{propositional variable}s~$\Obj p_i$
occurring in~$A$.

\begin{defn}
  The \emph{standard translation}~$\ST_x(A)$ is inductively defined
  as follows:
  \begin{enumerate}
  \tagitem{prvFalse}{\indcase{A}{\lfalse}{$\ST_x(\indfrm) = \lfalse$.}}{}
  \tagitem{prvTrue}{\indcase{A}{\lfalse}{$\ST_x(\indfrm) = \ltrue$.}}{}
  \item \indcase{A}{\Obj p_i}{$\ST_x(\indfrm) = \Atom{P_i}{x}$.}
  \tagitem{prvNot}{\indcase{A}{\lnot B}{$\ST_x(\indfrm) = \lnot \ST_x(B)$.}}{}
  \tagitem{prvAnd}{\indcase{A}{(B \land C)}{$\ST_x(\indfrm) = (\ST_x(B)
      \land \ST_x(C))$.}}{}
  \tagitem{prvOr}{\indcase{A}{(B \lor C)}{$\ST_x(\indfrm) = (\ST_x(B)
      \lor \ST_x(C))$.}}{}
  \tagitem{prvIf}{\indcase{A}{(B \lif C)}{$\ST_x(\indfrm) = (\ST_x(B)
      \lif \ST_x(C))$.}}{}
  \tagitem{prvIff}{\indcase{A}{(B \liff C)}{$\ST_x(\indfrm) = (\ST_x(B)
      \liff \ST_x(C))}$.}{}
  \tagitem{prvBox}{\indcase{A}{\Box B}{$\ST_x(\indfrm) =
      \lforall[y][(\Atom{Q}{x,y} \lif \ST_y(B))]$.}}{}
  \tagitem{prvDiamond}{\indcase{A}{\Diamond B}{$\ST_x(\indfrm) =
      \lexists[y][(\Atom{Q}{x,y} \land \ST_y(B))]$.}}{}
  \end{enumerate}
\end{defn}

For instance, $\ST_x(\Box p \lif p)$ is $\lforall[y][(\Atom{Q}{x,y}
  \lif \Atom{P}{y})] \lif \Atom{P}{x}$. Any \olup{structure} for the
language of~$\ST_x(A)$ requires a domain, a two-place relation
assigned to~$Q$, and subsets of the domain assigned to the one-place
\olup{predicate}s~$P_i$. In other words, the components of such
\olup{structure} are exactly those of a model for~$A$: the domain is
the set of worlds, the two-place relation assigned to~$Q$ is the
accessibility relation, and the subsets assigned to~$P_i$ are just the
assignments~$V(\Obj p_i)$. It won't surprise that satisfaction of $A$
in a modal model and of $\ST_x(A)$ in the corresponding \olup{structure}
agree:

\begin{prop}\ollabel{prop:st}
  Let $\mModel{M} = \tuple{W, R, V}$, $\Struct{M'}$ be the first-order
  \olup{structure} with $\Domain{M'} = W$, $\Assign{Q}{M'} = R$, and
  $\Assign{P_i}{M'} = V(\Obj p_i)$, and $s(x) = w$. Then
  \[
  \mSat{M}{A}[w] \text{ iff } \Sat{M'}{\ST_x(A)}[s]
  \]
\end{prop}

\begin{proof}
  By induction on~$A$.
\end{proof}

\begin{prop}
  Suppose $A$ is a modal \olup{formula} and $\mModel{F} = \tuple{W, R}$
  is a frame. Let $\Struct{F'}$ be the first-order \olup{structure} with
  $\Domain{F'} = W$ and $\Assign{Q}{F'} = R$, and let $A'$ be the
  second-order formula
  \[
  \lforall[X_1][\dots\lforall[X_n][\lforall[x][
        \SSubst{\ST_x(A)}{\subst{X_1}{P_1}, \dots,
          \subst{X_n}{P_n}}]]],
  \]
  where $P_1$, \dots, $P_n$ are all one-place \olup{predicate}s
  in~$\ST_x(A)$. Then
  \[
  \mModel{F} \Entails A \text{ iff } \Sat{F'}{A'}
  \]
\end{prop}

\begin{proof}
  $\Sat{F'}{A'}$ iff for every \olup{structure}~$\mModel{M'}$ where
  $\Assign{P_i}{M'} \subseteq W$ for $i = 1$, \dots,~$n$, and for
  every $s$ with $s(x) \in W$, $\Sat{M'}{\ST_x(A)}[s]$. By
  \olref{prop:st}, that is the case iff for all models~$\mModel{M}$
  based on~$\mModel{F}$ and every world~$w \in W$, $\mSat{M}{A}[w]$,
  i.e., $\mModel{F} \Entails A$.
\end{proof}

\begin{defn}
  A class~$\mClass{F}$ of frames is \emph{second-order definable} if
  there is \olup{sentence}~$A$ in the second-order language with a
  single two-place \olup{predicate}~$P$ and quantifiers only over monadic
  set variables such that $\mModel{F} = \tuple{W, R} \in \mClass{F}$
  iff $\Sat{M}{A}$ in the \olup{structure}~$\Struct{M}$ with
  $\Domain{M} = W$ and $\Assign{P}{M} = R$.
\end{defn}

\begin{cor}
  If a class of frames is definable by \olup{formula}~$A$, the
  corresponding class of accessibility relations is definable by a
  monadic second-order sentence.
\end{cor}

\begin{proof}
  The monadic second-order sentence~$A'$ of the preceding proof has
  the required property.
\end{proof}

As an example, consider again the \olup{formula}~$\Box p \lif p$. It
defines reflexivity. Reflexivity is of course first-order definable by
the \olup{sentence}~$\lforall[x][\Atom{Q}{x,x}]$. But it is also definable by the
monadic second-order \olup{sentence}
\[
\lforall[X][\lforall[x][(\lforall[y][(\Atom{Q}{x,y} \lif \Atom{X}{y})]
    \lif \Atom{X}{x})]].
\]
This means, of course, that the two sentences are equivalent. Here's
how you might convince yourself of this directly: First suppose the
second-order sentence is true in \olup{structure}~$\Struct{M}$. Since $x$ and $X$
are universally quantified, the remainder must hold for any $x \in W$
and set~$X \subseteq W$, e.g., the set $\Setabs{z}{Rxz}$ where $R =
\Assign{Q}{M}$. So, for any~$s$ with $s(x) \in W$ and $s(X) =
\Setabs{z}{Rxz}$ we have $\Sat{M}{\lforall[y][(\Atom{Q}{x,y} \lif
    \Atom{X}{y})] \lif \Atom{X}{x}}$. But by the way we've picked
$s(X)$ that means $\Sat{M}{\lforall[y][(\Atom{Q}{x,y} \lif
    \Atom{Q}{x,y})] \lif \Atom{Q}{x,x}}[s]$, which is equivalent to
$\Atom{Q}{x,x}$ since the antecedent is valid. Since $s(x)$ is
arbitrary, we have $\Sat{M}{\lforall[x][\Atom{Q}{x,x}]}$.

Now suppose that $\Sat{M}{\lforall[x][\Atom{Q}{x,x}]}$ and show that
$\Sat{M}{\lforall[X][\lforall[x][(\lforall[y][(\Atom{Q}{x,y} \lif
        \Atom{X}{y})] \lif \Atom{X}{x})]]}$. Pick any assignment~$s$,
and assume $\Sat{M}{\lforall[y][(\Atom{Q}{x,y} \lif
    \Atom{X}{y})]}[s]$. Let $s'$ be the $y$-variant of~$s$ with $s'(y)
= s(x)$; we have $\Sat{M}{\Atom{Q}{x,y} \lif \Atom{X}{y}}[s']$, i.e.,
$\Sat{M}{\Atom{Q}{x,x} \lif \Atom{X}{x}}[s]$. Since
$\Sat{M}{\lforall[x][\Atom{Q}{x,x}]}$, the antecedent is true, and we
have $\Sat{M}{\Atom{X}{x}}[s]$, which is what we needed to show.

Since some definable classes of frames are not first-order definable,
not every monadic second-order \olup{sentence} of the form~$A'$ is
equivalent to a first-order \olup{sentence}. There is no effective method
to decide which ones are.


\end{document}
